% Generated from locale/tr/content/computability/computability-theory/s-m-n.tex; do not hand-edit.


\olfileid[tr]{cmp}{thy}{smn}
\olsection{$s$-$m$-$n$ Teoremi}

\begin{explain}
Sıradaki teorem, nedeni birazdan açıklığa kavuşacak biçimde
``$s$-$m$-$n$ teoremi'' diye bilinir. Zor olan, teoremin tam olarak ne
söylediğini anlamaktır; ifadeyi anladıktan sonra sonuç oldukça açık görünür.
\end{explain}

\begin{thm}
\ollabel{thm:s-m-n}
  Her $n$ ve $m$ doğal sayı çifti için öyle bir ilkel özyinelemeli
  $s^m_n$ fonksiyonu vardır ki bütün
  $e,a_0,\dots,a_{m-1},y_0,\dots,y_{n-1}$ dizileri için
  \[
    \cfind{s^m_n(e,a_0,\dots,a_{m-1})}[n](y_0,\dots,y_{n-1})\simeq
    \cfind{e}[m+n](a_0,\dots,a_{m-1},y_0,\dots,y_{n-1})
  \]
  olur.
\end{thm}

\begin{explain}
$s^m_n$ fonksiyonunu \emph{programlar} üzerinde çalışan bir fonksiyon olarak
düşünmek yararlıdır. $s^m_n$, $(m+n)$-yerli bir fonksiyon için $e$ programı
ile sabitlenmiş $a_0,\dots,a_{m-1}$ girdilerini alır ve geriye kalan
argümanların $n$-yerli fonksiyonu için
$s^m_n(e,a_0,\dots,a_{m-1})$ programını döndürür.
\iftag{TMs}{$e$'yi bir Turing makinesinin betimi olarak düşünürseniz
$s^m_n(e,a_0,\dots,a_{m-1})$, $y_0,\dots,y_{n-1}$ girdisinde
$a_0,\dots,a_{m-1}$ öğelerini girdi dizisinin önüne ekleyen ve $e$'yi
çalıştıran Turing makinesidir. Her $s^m_n$ yalnızca uygun Turing makinesinin
kodunu bulan ilkel özyinelemeli bir fonksiyondur.}{}
\end{explain}

